\textbf{6.\ (Dormitory shopping).}
Anna and Belle live in the same dormitory. They go shopping every day, and
shopping takes \(20\) minutes. They decided to cooperate for \(4\) days:
on Monday Anna buys food for both of them, on Tuesday Belle does the same,
on Wednesday Anna again, and on Thursday Belle again.

Buying food for the mate takes \(5\) extra minutes. Thus, the assigned shopper
has an incentive to break the rule. If any agent does not buy food for the mate,
then both go shopping by themselves for the rest of the game.

Formalize the situation as a dynamic game and find all subgame perfect Nash
equilibria.

\textbf{Solution.}

Let
\[
c=\text{``buy for both''},\qquad d=\text{``buy only for yourself''}
\]

If someone chooses \(d\), cooperation ends and both shop separately for the rest
of the game. Payoffs are negative total shopping time.

Use backward induction.

On Thursday, Belle chooses between
\[
c:25,\qquad d:20
\]
Thus Belle chooses \(d\).

On Wednesday, Anna knows Belle will defect on Thursday. If Anna cooperates, her
remaining cost is
\[
25+20=45
\]
If Anna defects, her remaining cost is
\[
20+20=40
\]
Thus Anna chooses \(d\).

On Tuesday, Belle knows Anna will defect on Wednesday. If Belle cooperates, her
remaining cost is
\[
25+20+20=65.
\]
If Belle defects, her remaining cost is
\[
20+20+20=60.
\]
Thus Belle chooses \(d\).

On Monday, Anna knows Belle will defect on Tuesday. If Anna cooperates, her
total cost is
\[
25+20+20+20=85.
\]
If Anna defects, her total cost is
\[
20+20+20+20=80.
\]
Thus Anna chooses \(d\).

Therefore the unique subgame perfect Nash equilibrium is
\[
s_A=(d,d),\qquad s_B=(d,d)
\]

That is, each player defects whenever it is her turn to buy.

On the equilibrium path, Anna defects immediately on Monday. Then both shop
alone every day. Hence total shopping times are
\[
(80,80)
\]

The agreement is therefore not subgame perfect.
